\section{从削藩到轮台：帝国扩张如何改变中枢}
\label{sec:western-han-expansion-regency}

前一五四年春，晁错主持削减诸侯王封地，吴、楚等七国举兵。王国并非汉廷之外的陌生世界：同姓宗室握有封号、食邑、官属与地方关系，中央则以违法削地来重申自己的裁决权。景帝杀晁错没有使叛军退兵，反而表明争端已经超过一位大臣的存废。周亚夫在昌邑一线牵制、断其粮道，叛乱随后瓦解；兵力和战场细节在《汉书》《史记》之间仍有异文，但这场战争的结果很清楚：王国没有被废除，任官、司法与军政的自主空间却在此后继续收缩。

\subsection{把王国问题变成帝国能力}

削藩之后的中央并不是一部已经完备的机器。前一四一年，刘彻即位；礼制、人事和外戚之间的早期争论，显示年轻皇帝仍需在既有权力关系中寻找位置。前一三六年置五经博士、前一三四年令郡国举孝廉，分别把经典解释和地方评价接入官职体系。它们不能证明全国已按同一种课程培养官员，亦不能把“孝廉”直接等同于后来稳定的选官制度；不过，朝廷确实正在把能被书写、考核和推荐的人，变成可从郡国吸纳的行政资源。

边疆的选择也在改变中枢。张骞前一三九年出使，最初谋求联络大月氏的目标和行程细节仍有争议；他所带回的西方信息，却让朝廷的地理想象不再只围绕北方边郡。前一三三年马邑伏击未成，汉匈关系转入长期战争。此后卫青、霍去病等人的出击、河西走廊的经营与漠北战役，构成了汉朝从防守转向远征的年序。斩获、行军距离和“击破”究竟意味着控制到何处，不能照单全收；军队必须跨越草原、山地和补给线，才是这些叙事共同留下的约束。

\subsection{河西的道路与中枢的账簿}

前一二一年，霍去病两次出陇西，浑邪王随后率部降汉；前一二〇年至前一一一年，武威、酒泉、张掖、敦煌等郡陆续设立。把这些年表连在一起，并不等于河西从此成为边界清楚、人口均一的内地。降众安置、郡县设置、屯戍与邮传都需要长期投入，传世史书又主要从长安的视角记录这一过程。居延、肩水金关和悬泉置等边塞文书能够提示吏员、军粮和出入关的日常尺度，但它们是特定地点的档案，不能替河西全域或全帝国作证。

远征和经营需要钱粮。前一一九年前后，朝廷收铸五铢钱、强化盐铁专营；均输、平准又试图把各地物资调到中枢可以掌握的运输与价格网络。钱范、钱币窖藏和《汉书·食货志》使制度的存在可以互相参照，却不能证明每一项措施在各地同时、同样有效。更可靠的判断是，财政措施同边防不宜拆开理解：当军队、道路与新设郡县都要资源时，皇帝和负责财政者便试图把铸钱、盐铁和货物流转编入更集中的调度。

前一一二年至前一〇八年，汉军又先后征南越、入滇、灭卫满朝鲜。番禺、益州和朝鲜半岛的战事，不能被写成由长安一声命令便完成的直线扩张；南越内部的派系、地方首领的选择、郡数和实际控制的速度都各有争议。但这些行动确实把南海、西南和东北的道路、港口与郡县任命纳入朝廷的战略视野。帝国的空间变大，也使每一次征发、安置和防守的成本更难由关中一地吸收。

\subsection{巫蛊之后的收束}

前九一年，巫蛊案波及太子刘据，长安发生内战，太子败死。关于所谓巫蛊证据、告发者与责任分配，《汉书》〈武五子传〉和《史记》所保存的材料带有浓重的事后叙事；不必把宫廷传闻补成一次可以逐刻重演的密谋。可见的后果是，皇位继承与宫廷信任同时破裂，长期扩张所形成的中枢压力在都城内部爆发。

前八九年轮台诏拒绝继续在西域大规模屯田远征。诏书的存在可靠，却不能被夸写成武帝朝忽然放弃边疆，或一项政策立即抵销此前数十年的战争和财政重组。两年后武帝去世，少帝刘弗陵即位，霍光等受遗诏辅政。由远征、财政、官僚与继承共同塑成的国家，至此需要依靠辅政者维持其日常运转；后来的朝廷要处理的，已不只是如何扩大疆域，更是如何约束已被扩大的中枢。

\subsection{材料的视角}

这一组年序主要依靠《汉书》帝纪、列传和〈食货志〉，并与《史记》相关篇章、《盐铁论》、西域材料和边塞简牍对读。《汉书》能让人看见朝廷如何给削藩、专营和边疆行动排列次序，却也以后来形成的儒家正统解释武帝政治；《史记》的人物叙事和《盐铁论》的论辩都不能自动成为制度实施的纪录。出土文书和钱币、城址材料提供必要的物质校正，但其地域和保存边界同样不容越过。本节因此只把战争、设郡、财政措施和继承危机写为相互牵连的选择，不用未获交叉证明的数字、边界或人物动机填满空白。
